%! Author = chtompki
%! Date = 1/14/26
% Example LaTeX document for GP111 - note % sign indicates a comment
\documentclass[12pt]{amsart}

% Default margins are too wide all the way around. I reset them here
\setlength{\hoffset}{-.75in}
\setlength{\textwidth}{6.5in}
\setlength{\voffset}{-.5in}
\setlength{\textheight}{9.0in}

\usepackage{amsmath}
\usepackage{amssymb}
\usepackage{amsthm}
\usepackage{pgfplots}
\usepackage{enumerate}
\usepackage{tikz}
\usetikzlibrary{shadows,arrows,arrows.meta,shapes,positioning,calc}
\usepackage{setspace}
\usepackage{siunitx}
\usepackage{enumitem}

\theoremstyle{conjecture}
\newtheorem{conjecture}{Conjecture}

\theoremstyle{definition}
\newtheorem{definition}{Definition}

\theoremstyle{question}
\newtheorem{question}{Question}

\theoremstyle{theorem}
\newtheorem{theorem}{Theorem}

\newenvironment{solution}
{\begin{proof}[Solution]}
{\end{proof}}

\title{Graph Theory:\\Homework 1\\ Book Section 1.1.}
\author{Rob Tompkins}
\date{\today}

\begin{document}

    \maketitle

    % Problem 1
    \textbf{1.} Consider the following four families of graphs:
    \[
        A=\{\text{paths}\}, \qquad
        B=\{\text{cycles}\}, \qquad
        C=\{\text{complete graphs}\}, \qquad
        D=\{\text{bipartite graphs}\}.
    \]
    For each pair of these families, determine all isomorphism classes of
    graphs that belong to both families.

    \begin{solution}
        We consider each of the six possible pairs.

        \begin{enumerate}[label=(\alph*)]
            \item $A\cap B$: paths and cycles.

            A path has two vertices of degree $1$ unless it is a trivial
            path, while every vertex of a cycle has degree $2$. Therefore,
            no graph is both a path and a cycle. Hence
            \[
                A\cap B=\varnothing.
            \]

            \item $A\cap C$: paths and complete graphs.

            The graphs $K_1$ and $K_2$ are paths. If $n\geq 3$, every
            vertex of $K_n$ has degree $n-1\geq 2$, whereas a path has
            vertices of degree at most $2$ and has endpoints of degree $1$.
            Thus
            \[
                A\cap C=\{K_1,K_2\}.
            \]

            \item $A\cap D$: paths and bipartite graphs.

            Every path is bipartite. To see this, color the vertices
            alternately along the path. Adjacent vertices then receive
            different colors. Therefore
            \[
                A\cap D=A.
            \]

            \item $B\cap C$: cycles and complete graphs.

            The graph $K_3$ is a cycle of length $3$, namely $C_3$.
            For $n\geq 4$, every vertex of $K_n$ has degree at least $3$,
            while every vertex of a cycle has degree $2$. Hence
            \[
                B\cap C=\{K_3\}=\{C_3\}.
            \]

            \item $B\cap D$: cycles and bipartite graphs.

            A cycle is bipartite if and only if it has even length.
            Therefore
            \[
                B\cap D=\{C_{2k}:k\geq 2\}.
            \]

            \item $C\cap D$: complete graphs and bipartite graphs.

            The graphs $K_1$ and $K_2$ are bipartite. If $n\geq 3$,
            $K_n$ contains a triangle, and therefore is not bipartite.
            Hence
            \[
                C\cap D=\{K_1,K_2\}.
            \]
        \end{enumerate}
    \end{solution}

    % Problem 2
    \textbf{2.} Prove that removing two opposite corner squares from an
    $8$-by-$8$ checkerboard leaves a subboard that cannot be partitioned
    into $1$-by-$2$ and $2$-by-$1$ rectangles. Using the same argument,
    make a general statement about all bipartite graphs.

    \medskip

    \noindent
    \textit{Hint:} An $8$-by-$8$ checkerboard can be represented by a graph.
    Each square corresponds to a vertex, and two vertices are adjacent if
    and only if the corresponding squares share a side.

    \begin{proof}
        Color the checkerboard in the usual way with alternating black and
        white squares. There are initially $32$ black squares and $32$
        white squares.

        Opposite corners of an $8$-by-$8$ checkerboard have the same color.
        Thus, after removing two opposite corner squares, the remaining
        board contains either
        \[
            30\text{ black squares and }32\text{ white squares},
        \]
        or
        \[
            32\text{ black squares and }30\text{ white squares}.
        \]

        Every $1$-by-$2$ or $2$-by-$1$ rectangle covers two adjacent
        squares. Since adjacent squares have opposite colors, every such
        rectangle covers exactly one black square and one white square.

        Therefore, any board that can be partitioned into these rectangles
        must contain the same number of black and white squares. The board
        obtained by deleting two opposite corners does not, so such a
        partition is impossible.

        In graph-theoretic language, let $G$ be a bipartite graph with
        partite sets $X$ and $Y$. A collection of edges covering every
        vertex exactly once is a perfect matching. Every edge in a
        bipartite graph contains one vertex from $X$ and one vertex from
        $Y$. Hence, if $G$ has a perfect matching, then necessarily
        \[
            |X|=|Y|.
        \]

        Thus, a bipartite graph whose two partite sets have different
        cardinalities cannot have a perfect matching.
    \end{proof}

    % Problem 3
    \textbf{3.} Let $G$ be a simple bipartite graph with $n$ vertices and $m$
    edges. Prove that
    \[
        m \leq \frac{n^2}{4}.
    \]

    \medskip

    \noindent
    \textit{Hint:} Suppose $X$ and $Y$ are partite sets of $G$ with
    $|X|=r$ and $|Y|=s$. First show that
    \[
        m \leq rs,
    \]
    and use it to prove
    \[
        m \leq \frac{n^2}{4}.
    \]

    \begin{proof}
        Let $X$ and $Y$ be the two partite sets of $G$, with
        \[
            |X|=r
            \qquad\text{and}\qquad
            |Y|=s.
        \]
        Since every edge of a bipartite graph joins a vertex of $X$ to a
        vertex of $Y$, and since $G$ is simple, there can be at most one
        edge for each pair $(x,y)\in X\times Y$. Therefore
        \[
            m\leq rs.
        \]

        Since
        \[
            r+s=n,
        \]
        we have
        \[
            (r-s)^2\geq 0.
        \]
        Expanding gives
        \[
            r^2-2rs+s^2\geq 0.
        \]
        Hence
        \[
            r^2+2rs+s^2\geq 4rs,
        \]
        so
        \[
            (r+s)^2\geq 4rs.
        \]
        Since $r+s=n$,
        \[
            n^2\geq 4rs,
        \]
        and therefore
        \[
            rs\leq \frac{n^2}{4}.
        \]

        Combining the two inequalities gives
        \[
            \boxed{m\leq \frac{n^2}{4}}.
        \]

        Equality occurs when the two partite sets are as equal in size as
        possible and every possible edge between them is present.
    \end{proof}

    % Problem 4
    \textbf{4.} Let $G$ be a graph with girth $4$ in which every vertex has
    degree $k$. Prove that $G$ has at least $2k$ vertices. Determine all
    such graphs with exactly $2k$ vertices.

    \begin{proof}
        Since $G$ has girth $4$, it contains no triangles.

        Choose a vertex $v\in V(G)$. Since every vertex has degree $k$,
        $v$ has exactly $k$ neighbors. Let
        \[
            N(v)=\{v_1,v_2,\ldots,v_k\}.
        \]

        Because $G$ contains no triangles, no two vertices in $N(v)$ can
        be adjacent. Otherwise, if $v_i$ and $v_j$ were adjacent, then
        \[
            v,v_i,v_j,v
        \]
        would form a $3$-cycle.

        Now choose $v_1\in N(v)$. Since $\deg(v_1)=k$, the vertex $v_1$
        has $k-1$ neighbors other than $v$. None of these $k-1$ vertices
        can lie in $N(v)$, since that would create a triangle containing
        $v$.

        Thus the graph contains at least
        \[
            1+k+(k-1)=2k
        \]
        vertices. Hence
        \[
            |V(G)|\geq 2k.
        \]

        Now suppose $|V(G)|=2k$. Define
        \[
            X=N(v)
        \]
        and
        \[
            Y=V(G)\setminus X.
        \]
        Since $|X|=k$ and $|V(G)|=2k$, we have $|Y|=k$.

        There are no edges between vertices of $X$, since such an edge
        would create a triangle with $v$. Each vertex of $X$ has degree
        $k$, and there are exactly $k$ vertices in $Y$. Therefore every
        vertex of $X$ must be adjacent to every vertex of $Y$.

        Consequently,
        \[
            G\cong K_{k,k}.
        \]

        Conversely, for $k\geq 2$, the complete bipartite graph $K_{k,k}$
        is $k$-regular, has $2k$ vertices, and has girth $4$. Therefore
        the graphs attaining equality are precisely
        \[
            \boxed{K_{k,k}}.
        \]
    \end{proof}

    % Problem 5
    \textbf{5.} Let $P$ be the Petersen graph. Use the definition of $P$ given
    in terms of $2$-element subsets of a set of size $5$ to prove the
    following.

    \begin{enumerate}[label=(\alph*)]
        \item $P$ contains a cycle of length $6$.

        \item $P$ contains a cycle of length $8$.

        \item $P$ does not contain a cycle of length $7$.

        \medskip

        \noindent
        \textit{Hint:} For contradiction, assume that $C$ is a cycle in
        $P$ of length $7$. First, explain why every edge of $P$ whose
        endpoints both lie in $C$ must be an edge of $C$.

        Let $u$ be a vertex of $C$ and let $v,w$ be vertices of $C$
        farthest from $u$. Show that there is some vertex
        \[
            x \notin V(C)
        \]
        such that $u$ and $v$ are both adjacent to $x$, and $u$ and $w$
        are both adjacent to $x$. Conclude that $\{w,v,x\}$ are the
        vertices of a $3$-cycle in $P$, which contradicts the fact that
        $P$ is triangle-free.
    \end{enumerate}

    \begin{solution}
        Recall that the vertices of the Petersen graph are the $2$-element
        subsets of
        \[
            \{1,2,3,4,5\},
        \]
        and two vertices are adjacent precisely when the corresponding
        $2$-element subsets are disjoint.

        \medskip

        \noindent
        \textbf{(a)} Consider the sequence
        \[
            12,\;34,\;15,\;23,\;45,\;13,\;12.
        \]
        Consecutive $2$-element subsets are disjoint:
        \[
            \begin{aligned}
                12\cap34&=\varnothing,\\
                34\cap15&=\varnothing,\\
                15\cap23&=\varnothing,\\
                23\cap45&=\varnothing,\\
                45\cap13&=\varnothing,\\
                13\cap12&=\varnothing.
            \end{aligned}
        \]
        Thus these six distinct vertices form the cycle
        \[
            12-34-15-23-45-13-12.
        \]
        Hence $P$ contains a cycle of length $6$.

        \medskip

        \noindent
        \textbf{(b)} Consider
        \[
            12,\;34,\;15,\;23,\;14,\;25,\;13,\;45,\;12.
        \]
        Again, every pair of consecutive $2$-element subsets is disjoint:
        \[
            \begin{aligned}
                12\cap34&=\varnothing,\\
                34\cap15&=\varnothing,\\
                15\cap23&=\varnothing,\\
                23\cap14&=\varnothing,\\
                14\cap25&=\varnothing,\\
                25\cap13&=\varnothing,\\
                13\cap45&=\varnothing,\\
                45\cap12&=\varnothing.
            \end{aligned}
        \]
        Therefore
        \[
            12-34-15-23-14-25-13-45-12
        \]
        is an $8$-cycle in $P$.

        \medskip

        \noindent
        \textbf{(c)} Suppose, for contradiction, that $P$ contains a
        $7$-cycle $C$.

        First observe that $C$ cannot have a chord. Indeed, a chord of a
        $7$-cycle divides the cycle into two smaller cycles whose lengths
        sum to $9$. Since the Petersen graph has girth $5$, neither of
        these cycles can have length less than $5$. This is impossible
        because two integers each at least $5$ cannot sum to $9$.
        Therefore every edge of $P$ whose endpoints lie in $C$ must
        already be an edge of $C$.

        Since the Petersen graph is $3$-regular, each vertex of $C$ has
        exactly two neighbors on $C$ and exactly one neighbor outside
        $C$.

        There are $10$ vertices in the Petersen graph, so exactly three
        vertices lie outside $C$. Let these vertices be
        \[
            x,y,z.
        \]

        The seven vertices of $C$ each have exactly one edge leading
        outside $C$. Thus there are seven edges from $C$ to
        $\{x,y,z\}$.

        Since $x,y,z$ each have degree $3$, the sum of their degrees is
        \[
            3+3+3=9.
        \]
        If $e$ denotes the number of edges among $x,y,z$, then those
        internal edges contribute $2e$ to this degree sum. Hence
        \[
            7+2e=9,
        \]
        so
        \[
            e=1.
        \]

        Thus among the three vertices outside $C$, exactly one pair is
        adjacent.

        Let $u$ be a vertex of $C$, and let $v$ and $w$ be the two
        vertices of $C$ at distance $3$ from $u$ along the cycle. Since
        the Petersen graph has diameter $2$, both $u$ and $v$ must have
        a common neighbor, and both $u$ and $w$ must have a common
        neighbor.

        These common neighbors cannot lie on $C$, since $C$ has no
        chords. Therefore the common neighbors must lie outside $C$.

        But $u$ has only one neighbor outside $C$. Call it $x$. Hence
        $x$ must be a common neighbor of $u$ and $v$, and also a common
        neighbor of $u$ and $w$. Thus
        \[
            xv,\;xw\in E(P).
        \]

        Since $v$ and $w$ are adjacent on the $7$-cycle $C$, we also have
        \[
            vw\in E(P).
        \]

        Therefore the three vertices
        \[
            v,\;w,\;x
        \]
        form a triangle
        \[
            v-x-w-v.
        \]

        This contradicts the fact that the Petersen graph is
        triangle-free. Therefore the Petersen graph cannot contain a
        cycle of length $7$.
    \end{solution}

\end{document}